\section{Introduction}

Deploying large neural networks on resource-constrained sensors and wearables
demands aggressive model compression without sacrificing task accuracy.
Structured pruning \cite{han2015deep} and quantisation \cite{jacob2018quantization}
are well-studied, yet few approaches simultaneously control gradient dynamics
during pruning, leading to premature weight death or unstable fine-tuning.

Recent lottery-ticket analyses \cite{frankle2019lottery} suggest that sparse
subnetworks capable of matching full-model accuracy exist within randomly
initialised networks; however, reliably identifying these subnetworks during
training remains an open problem, especially under strict energy budgets.

\textbf{Contributions.}
\begin{itemize}[leftmargin=1.2em,topsep=2pt,itemsep=1pt]
  \item A differentiable threshold gate $\sigma_\tau$ that learns per-layer
        sparsity targets end-to-end.
  \item A gradient rescaling rule that prevents gradient vanishing in pruned
        channels while preserving sparsity pressure.
  \item State-of-the-art accuracy–energy trade-off on five MCU platforms.
\end{itemize}

% ─── 2. Method ───────────────────────────────────────────────────────────────
